% GENERATED FILE. Edit canonical lessons, then run npm run generate:books.
% canonical-source-hash: fnv1a64:1cf00301618f8ea9
% canonical-lessons: SA-C185-kiyat, SA-C185-kidrsah, SA-C185-hi, SA-C185-nunam, SA-C185-avasyam

\chapter{How much?, Of what kind?, For, Surely, Certainly}
\label{ch:sa-how-much-of-what-kind-for-surely-certainly}
\hlchaptermodality{185}

\begin{chapteropening}
\textbf{By the end of this chapter:} \emph{I can say how much?, of what kind?, for, surely and certainly.}

Complete the last lesson of chapter 185: I can say how much?, of what kind?, for, surely and certainly.
\end{chapteropening}

\section[\texorpdfstring{\sk{कियत्} (kiyat)}{कियत् (kiyat)}]{\sk{कियत्} (kiyat) — how much?}

\label{lesson:SA-C185-kiyat}



\begin{quote}
Before the new one: say the Sanskrit for loudly, then the Sanskrit for softly.
\end{quote}

\subsection*{You'll want to know: \sk{कियत्}}
\textbf{\sk{कियत्}} — \emph{kiyat} — \textquotedblleft{}how much?\textquotedblright{}.

\begin{grammarlens}[title={What you've built}]
One more word for this chapter.
\end{grammarlens}

\subsection*{Your turn}
\begin{itemize}
\raggedright
  \item \emph{Say it:} \emph{kiyat}
  \item \emph{Say it:} \emph{kiyat}, once more
  \item \emph{Say it:} say \emph{kiyat}
  \item \emph{Recall:} say the Sanskrit for loudly, then the Sanskrit for softly, then say \emph{kiyat} again
\end{itemize}

\begin{tcolorbox}[breakable,colback=teal!4,colframe=teal!35!black,title={Before you move on}]
What does \sk{कियत्} mean? (\textquotedblleft{}How much?\textquotedblright{}.) Say it once more.
\end{tcolorbox}
\section[\texorpdfstring{\sk{कीदृशः} (kīdṛśaḥ)}{कीदृशः (kīdṛśaḥ)}]{\sk{कीदृशः} (kīdṛśaḥ) — of what kind?}

\label{lesson:SA-C185-kidrsah}



\begin{quote}
Before the new one: say the Sanskrit for softly, then the Sanskrit for how much?.
\end{quote}

\subsection*{You'll want to know: \sk{कीदृशः}}
\textbf{\sk{कीदृशः}} — \emph{kīdṛśaḥ} — \textquotedblleft{}of what kind?\textquotedblright{}.

\begin{grammarlens}[title={What you've built}]
One more word for this chapter.
\end{grammarlens}

\subsection*{Your turn}
\begin{itemize}
\raggedright
  \item \emph{Say it:} \emph{kīdṛśaḥ}
  \item \emph{Say it:} \emph{kīdṛśaḥ}, once more
  \item \emph{Say it:} say \emph{kīdṛśaḥ}
  \item \emph{Recall:} say the Sanskrit for softly, then the Sanskrit for how much?, then say \emph{kīdṛśaḥ} again
\end{itemize}

\begin{tcolorbox}[breakable,colback=teal!4,colframe=teal!35!black,title={Before you move on}]
What does \sk{कीदृशः} mean? (\textquotedblleft{}Of what kind?\textquotedblright{}.) Say it once more.
\end{tcolorbox}
\section[\texorpdfstring{\sk{हि} (hi)}{हि (hi)}]{\sk{हि} (hi) — for, indeed}

\label{lesson:SA-C185-hi}



\begin{quote}
Before the new one: say the Sanskrit for how much?, then the Sanskrit for of what kind?.
\end{quote}

\subsection*{You'll want to know: \sk{हि}}
\textbf{\sk{हि}} — \emph{hi} — \textquotedblleft{}for, indeed\textquotedblright{}.

\begin{grammarlens}[title={What you've built}]
One more word for this chapter.
\end{grammarlens}

\subsection*{Your turn}
\begin{itemize}
\raggedright
  \item \emph{Say it:} \emph{hi}
  \item \emph{Say it:} \emph{hi}, once more
  \item \emph{Say it:} say \emph{hi}
  \item \emph{Recall:} say the Sanskrit for how much?, then the Sanskrit for of what kind?, then say \emph{hi} again
\end{itemize}

\begin{tcolorbox}[breakable,colback=teal!4,colframe=teal!35!black,title={Before you move on}]
What does \sk{हि} mean? (\textquotedblleft{}For, indeed\textquotedblright{}.) Say it once more.
\end{tcolorbox}
\section[\texorpdfstring{\sk{नूनम्} (nūnam)}{नूनम् (nūnam)}]{\sk{नूनम्} (nūnam) — surely}

\label{lesson:SA-C185-nunam}



\begin{quote}
Before the new one: say the Sanskrit for of what kind?, then the Sanskrit for for.
\end{quote}

\subsection*{You'll want to know: \sk{नूनम्}}
\textbf{\sk{नूनम्}} — \emph{nūnam} — \textquotedblleft{}surely\textquotedblright{}.

\begin{grammarlens}[title={What you've built}]
One more word for this chapter.
\end{grammarlens}

\subsection*{Your turn}
\begin{itemize}
\raggedright
  \item \emph{Say it:} \emph{nūnam}
  \item \emph{Say it:} \emph{nūnam}, once more
  \item \emph{Say it:} say \emph{nūnam}
  \item \emph{Recall:} say the Sanskrit for of what kind?, then the Sanskrit for for, then say \emph{nūnam} again
\end{itemize}

\begin{tcolorbox}[breakable,colback=teal!4,colframe=teal!35!black,title={Before you move on}]
What does \sk{नूनम्} mean? (\textquotedblleft{}Surely\textquotedblright{}.) Say it once more.
\end{tcolorbox}
\section[\texorpdfstring{\sk{अवश्यम्} (avaśyam)}{अवश्यम् (avaśyam)}]{\sk{अवश्यम्} (avaśyam) — certainly, necessarily}

\label{lesson:SA-C185-avasyam}



\begin{quote}
Before the new one: say the Sanskrit for for, then the Sanskrit for surely.
\end{quote}

\subsection*{You'll want to know: \sk{अवश्यम्}}
\textbf{\sk{अवश्यम्}} — \emph{avaśyam} — \textquotedblleft{}certainly, necessarily\textquotedblright{}.

\begin{grammarlens}[title={What you've built}]
One more word for this chapter.
\end{grammarlens}

\subsection*{Your turn}
\begin{itemize}
\raggedright
  \item \emph{Say it:} \emph{avaśyam}
  \item \emph{Say it:} \emph{avaśyam}, once more
  \item \emph{Say it:} say \emph{avaśyam}
  \item \emph{Recall:} say the Sanskrit for for, then the Sanskrit for surely, then say \emph{avaśyam} again
\end{itemize}

\begin{tcolorbox}[breakable,colback=teal!4,colframe=teal!35!black,title={Before you move on}]
What does \sk{अवश्यम्} mean? (\textquotedblleft{}Certainly, necessarily\textquotedblright{}.) Say it once more.
\end{tcolorbox}
